% auto-generated by make exhibits -- do not edit
\begin{table}[!htbp]\centering\footnotesize
\caption{Institutional and demographic interactions with the rule (appendix)}\label{tab:r3}
\setlength{\tabcolsep}{3.5pt}
\begin{tabular}{@{}lccccc@{}}
\toprule
 & (1) & (2) & (3) & (4) & (5) \\
 & Voted \$ share & GO share & NC share & Voted \$ share & Voted \$ share \\
 & TEL interaction & TEL interaction & TEL interaction & $\times$ homeownership & $\times$ county Dem share \\
\midrule
Strict rule (main effect) & +0.205 & -0.109 & +0.005 & +0.096 & +0.194 \\ \addlinespace
Interaction & -0.185 & -1.227 & -0.638 & -0.047 & -0.651\sig{**} \\
 & (0.338) & (0.821) & (0.494) & (0.395) & (0.326) \\ \addlinespace
\midrule
Observations & 396 & 401 & 358 & 13,928 & 13,928 \\
State clusters & 34 & 34 & 34 & 48 & 48 \\
\bottomrule
\end{tabular}
\begin{minipage}{0.96\linewidth}\vspace{2pt}\scriptsize
\emph{Notes:} Columns (1)--(3): the roughly 570-city subpanel with tax-and-expenditure-limit (TEL) stringency; the interaction asks whether TELs sharpen or blunt the rule; the subpanel is the exit-rich class where weak rule effects are predicted, so power is limited by design. Columns (4)--(5): national entity panel, strict rule interacted with the centred moderator; these moderate the first-stage channel, not the RD issuance effect (the causal moderator tests are in Section 8). Entity panel (90,604 local governments; estimation samples below are the corpus-active units for each outcome). Weighted least squares; Census-region fixed effects; controls: ln population (ln total revenue for special districts), homeownership, share 65+, racial fractionalisation, ln median household income, county Democratic presidential share 2020. Standard errors clustered by state in parentheses. The strict-rule indicator is the PRELIMINARY pass-1 coding (verification pass in progress), so coefficients read as first-stage associations, not causal estimates. \sig{*} $p<0.1$, \sig{**} $p<0.05$, \sig{***} $p<0.01$.
\end{minipage}
\end{table}
